% chapters/ch16.tex —— 第 16 章
% 体例：具体失败 → 设计问题 → 机制解释 → 对照实验 → 条件变化 → 选择原则
% 数字身份纪律：
%   - 本书实测（E01 / E02 / E03 三次运行）一律 \measured{}；
%   - 外部论文报告值（2607.07663 / 2607.12227 / 2602.19320）一律 \cited{}；
%   - 由上述数字算出的倍数、比例与合计用 \calcnum{}。
% 判决纪律：分别核对设计、范围、精度与方向；不以跨零或单个样本数门槛自动判决。
\chapter{理论边界与研究议程}
\label{ch16}

\chaptergoal{已有试验与实现审计能支持怎样的持续改进主张？}
\chapterprops{本章把机制与案例归纳为P1、P1b、P2、P2b、P3、P3b、P4、P5、P6、P7十个研究问题，分别连接已有观察、适用范围与后续设计；P8作为扩展方向另列。}
\chapterfig{图~\ref{fig:ch16-scope}：从有限比较到更强主张，需要扩展哪些评价}
\chapterpitfall{一个点估计 \measured{+0.0500}、区间 \measured{$[-0.1000, +0.1750]$} 的结果，
被写成“持久更新有效”}

% ---------- 1. 具体失败 ----------
\begin{failurebox}
本书的 E01、E02 留存了运行轨迹与汇总，E03 的十臂已于 2026-09-28 全部跑完
（最后一臂 \texttt{fit\_domain\_cross} 由跨日期的 \texttt{--resume-latest} 续跑补齐）。
这些运行提供了读数，也暴露了尚未满足的比较条件。
自进化研究需要把不同强度的问题写成不同的比较：预算分配、经验迁移与长期保持，各有自己的对象和范围。

E01 在 \measured{16} 道封存题、每题 \measured{5} 次运行上，
格 4 相对格 1 的差为 \measured{+0.0500}，
95\% 自助区间为 \measured{$[-0.1000,+0.1750]$}。
区间兼容两种方向；实现审计还确认格1并非重试后选择，完整成本也未对齐；
这组读数描述了已实现的四种方案，并把研究问题P1引向匹配重试策略和完整预算的比较。

E02五个更新臂的点估计低于其无记忆基线。E03拆分A按第5章的记录采用边界阅读，其中记忆臂的直接比较仍可复核。负向点估计与记录冲突需要分别处理。
本章将读数、设计有效性、估计精度和命题范围逐项对应，明确下一条证据应补在哪里。
\end{failurebox}

\designquestion{前面十五章分别处理了机制、检验、成本与防线。
本章只回答一个问题：\textbf{把这些结果放在一起，哪些主张是可以往外说的，哪些不是}。
本章将已测边界对应到具体运行，将未测问题标为理论条件或研究设计，不用数字装饰尚无观察的部分。}

第 1 章定下的三条公理里，公理 1 说验证器 $\mathcal V$ 默认冻结。
这条公理在第 3 章至第 14 章一直是默认前提，本章要说明它的代价：
\textbf{冻结 $\mathcal V$ 提供了共同参照，但参照是否可靠、是否独立、是否覆盖所问问题仍须另行检查}。
本章先核对泛化范围、学习条件与反馈，再区分长期稳定性、闭环风险与目标演化。十个研究问题把这些机制认识连接到具体比较；后续设计分别改善识别、精度与覆盖。

% ---------- 2. 有限评测与未知任务 ----------
\section{有限评测与未知任务上的泛化}
\label{sec:ch16-eval}

\begin{sloppypar}
自进化的迁移主张最终要回答：\textbf{这一次学到的东西，下一次在别处还有用吗？}
本节先区分三种“别处”。它们改变的条件不同，失效机制可能重叠，需要分别报告。
\end{sloppypar}

\begin{table}[htbp]
\centering
\caption{三种“别处”：各自的失败方式与本书已有的实测状态}
\label{tab:ch16-eval}
\begin{footnotesize}
\begin{tabular}{>{\raggedright\arraybackslash}p{2.4cm}>{\raggedright\arraybackslash}p{3.6cm}>{\raggedright\arraybackslash}p{4.0cm}>{\raggedright\arraybackslash}p{3.4cm}}
\toprule
泛化到 & 具体含义 & 本书实测 & 状态 \\
\midrule
同族新题 & 同一任务族内没做过的题 & E01整体格4减格1为\measured{+0.0500}（\measured{$[-0.1000, +0.1750]$}），但封存集混合实际训练已见与未见族 & 整体读数不能单独代表本行 \\
\midrule
新任务族 & 换一类失败模式 & E03拆分A：记忆臂直接比较可查，采用边界见第5章 & 已测方案与有限任务覆盖 \\
\midrule
新读取者 & 换基础模型或换运行脚手架 & E03 拆分 B：\texttt{base\_same} \measured{0.6875} 对 \texttt{base\_cross} \measured{0.2625} & 替代读取者的无记忆基线较低 \\
\bottomrule
\end{tabular}
\end{footnotesize}
\end{table}

\begin{sloppypar}
第一行须按实际训练暴露界定“同族”。E01任务划分为八族，但历史训练只执行了八道题、覆盖四族，封存集则含八族。表中的整体差值没有把这两类封存题分开，因此既不是纯同族迁移估计，也不是纯跨族迁移估计；目录中存在某族的训练题不等于该次运行见过它。
\end{sloppypar}

\begin{sloppypar}
第三种比较需要把基线差与记忆增量拆开。表中的 \texttt{base\_same} 和 \texttt{base\_cross} 都是不带新增记忆的基线，其差本身不是经验迁移损失。这里能说的是：替代模型在本轮特定任务、接口与预算下的基础表现较低；不能据此把它概括为所有任务上更弱。要估计经验迁移，须在每个读取者内部先比较记忆臂与对应基线，再比较这些增量。跨日期续跑和模型别名变化仍会限制解释。
\end{sloppypar}

\begin{sloppypar}
外部文献也报告过主干替换比较：两种记忆架构分别在所报任务范围内更换主干，答题分从 \cited{0.289} 掉到 \cited{0.102}（\texttt{SimpleMem}）、
从 \cited{0.781} 掉到 \cited{0.447}（\texttt{Nemori}），
可恢复格式偏差比例在这两种系统的比较中从 \cited{1.20\%} 升到 \cited{4.82\%}、
从 \cited{17.91\%} 升到 \cited{30.38\%}（2602.19320 表 4；该比例不是不可恢复记忆损坏率）\cite{anatomy-memory-2602}。
\textbf{这些数字全部是原论文报告值，不是本书实测}；它们说明的事情是清楚的：
更换主干可以同时改变任务表现与接口错误。这里的主干同时参与答题和记忆更新、整合，读写条件一起变化。这个联合变化不能分解出某个固定的“主干份额”，也不能排除主干与记忆机制的交互；本书据此要求配套静态基线。
\end{sloppypar}

\begin{pitfallbox}
把“在封存集上提升了”直接读成“泛化了”，是本章要挡住的第一个错误。
恰当的封存可以限制研究过程中使用评价题，却不能自动排除基础模型预训练暴露、模板相似或领域熟悉。历史实验的读取边界还存在已记录的缺口。三种“别处”必须分开报告，
否则 P1 与 P3b 会共用同一个数字而指向相反的结论。
\end{pitfallbox}

% ---------- 3. 探索、学习与系统选择 ----------
\section{探索、学习与系统选择}
\label{sec:ch16-select}

\begin{sloppypar}
第 2 章把“改什么”写成搜索空间 $\mathcal S$、把“怎么选”写成优化算法 $H$。
这个分解在本书里一直是工程性的：它描述一次更新长什么样。
要问它能不能一直做下去，需要的是条件，不是结构。
\end{sloppypar}

\begin{sloppypar}
综述\cite{rsi-survey-2607}在§7.1转述苏格拉底学习立场文章，提出反馈、覆盖与容量三个条件。本书借用它们组织工程问题，不把转述当作已经核验原始证明的定理，也不称为全领域共识。反馈是否有信息、经验覆盖是否充分、更新对象是否具有表示能力，需要分别操作化。“有信息”不等于信号来自另一个模型；“覆盖广”不等于题目数量多；“容量够”也不能只用参数个数回答。下表给出可检查的风险对应，而非一一对应的诊断法。

\end{sloppypar}

\begin{table}[htbp]
\centering
\caption{三个条件对应的风险与候选观测量}
\label{tab:ch16-cond}
\begin{footnotesize}
\begin{tabular}{>{\raggedright\arraybackslash}p{1.6cm}>{\raggedright\arraybackslash}p{4.0cm}>{\raggedright\arraybackslash}p{3.6cm}>{\raggedright\arraybackslash}p{4.2cm}}
\toprule
条件 & 可能出现的现象 & 讨论位置 & 本书的观测量 \\
\midrule
（a）反馈对齐 & 置信度高的错误被反复强化 & 第 4 章；E02 的噪声臂 & 反馈源分级；\texttt{mismatch\_rate}；接地频率 \\
\midrule
（b）经验覆盖 & 任务分布在共进化中收窄 & 第 13 章；第\ref{ch14}章\S\ref{sec:ch14-collapse} & 任务分布与族覆盖；按族、题分别报告 \\
\midrule
（c）容量 & 有限表示空间造成覆盖与干扰权衡 & 第 8 章 P4；第 8 章成本模型 & 覆盖、干扰与存储规模；成本另见E05 \\
\bottomrule
\end{tabular}
\end{footnotesize}
\end{table}

\begin{sloppypar}
内部修订、经验收益和无界增长是三个不同层面的主张。一次程序修改可以改善有限任务成绩，却没有因此证明持续迭代会永远获得正增量。反过来，某个形式模型对计算能力层级的限制，也不能直接推成现实系统无法改善速度、可靠性或有限预算内成绩。综述§7.1汇集了多种理论与经济论证；在没有逐一重建其对象、假设和结论前，本书不把它们合并成“自修改不可能提高能力”的通用定理。本章的研究问题关注的都是具有限定任务、资源和处理方式的比较。

\end{sloppypar}

\begin{sloppypar}
P7问修改器的改进能否迁移到不同任务族。HyperAgents v1迁出整套任务与元实现，目标阶段固定元实现并评价搜索结果\cite{hyperagents-2603}；v2又加入只迁任务、只迁元层及二者交叉的对照。在其数学证明评分配置中，只迁元Agent获得了接近整包迁移的中位增益\cite{hyperagents-2603-v2}。这把证据从整体迁移推进到组件作用，仍未逐一隔离元Agent内部的记忆、记录和生成机制。它为P7提供正向外部证据；本书E10的任务族、对照与资源范围仍须独立落实。

\end{sloppypar}

% ---------- 4. 非平稳环境 ----------
\subsection{反馈质量与验证能力的瓶颈}
\label{sec:ch16-feedback}

\begin{sloppypar}
E02 的结果在本章里是最重要的一个负结果：五个记忆臂的回合通过率点估计 \textbf{全部低于}无记忆基线 \measured{0.8125}。
自评臂最低 \measured{0.6125}（收益 \measured{$-0.2000$}），噪声臂 \measured{0.6625}（\measured{$-0.1500$}）。
历史题级区间中，自评减基线不跨零，但这不是完整层级的检验。
对第4章“自评减执行”与“测试反馈减自评”两项比较，以各自差值标准差反推三个百分点的区间半宽，分别得到\calcnum{592}与\calcnum{433}题的规划参考；它们不是上句自评、噪声各自减基线的规划值，
不是预注册样本量，也不是检验功效保证。
\end{sloppypar}

\begin{sloppypar}
但即使样本够，E02 也不足以判 P2b。原因在第 4 章已经说明：任务集 v1 的验收器是单元测试，
不是定理证明器；评审与自评又是同一后端、新上下文下的两种提示。测试文件可读，噪声臂只翻转标题而保留原始输出。
因此限制同时涉及操作定义、信息边界与归因，不只是样本量。
\textbf{这不是一个可以用更多样本修补的缺陷，它也影响当前处理究竟比较了什么，而不只限制外推范围。}
\end{sloppypar}

\begin{sloppypar}
如果系统连评审器也通过自己的输出更新，评审训练数据、验收数据和生成器输出之间的依赖关系就成了新增对象。改进评审器可以有用，但它的自评分上升不能单独说明误判减少。一个可检查的比较是冻结一组独立标注或可执行的任务，分别测量新旧评审器的误接纳与误拒绝，再检查这些变化怎样影响后续选择。这是本书提出的评价设计，不是已经执行的验证器训练实验。

\end{sloppypar}

\begin{sloppypar}
模型评审器可能错误，也可能在工具、分解或多次核查的帮助下评价超过一次直接回答能力的产物。因此，不宜把某个单模型的做题成绩直接当成所有评审方案的数值上限。操作上应报告评审校准样本、错误类型、外部核查覆盖和未核查部分；没有独立核查就明确缺口。封存题目主要约束信息使用，不能自动保证评审可靠，也不能证明基础模型预训练时从未见过相似材料。

\end{sloppypar}

\begin{principlebox}
综述\cite{rsi-survey-2607}所选的\cited{1,250}篇论文中，基础、限制与安全家族有\cited{60}篇，按该语料分母计算为\calcnum{4.8\%}。这是该检索与分类方案中的比例，不是所有研究的完整普查，也不能换算为已经投入的工时或资金。本书借它提出测量问题：对于一次持续更新，是否有足够记录判断改善了什么、牺牲了什么、哪些部分未被检验？统一预算和证据索引服务于这个问题，它们本身不增加算法成绩。

\end{principlebox}

% ---------- 6. 元进化与开放式 ----------
\section{非平稳环境中的稳定性、遗忘与适应}
\label{sec:ch16-nonstationary}

\begin{sloppypar}
第11章引用的Self-Challenging报告过GRPO训练不充分调参时的突然崩溃，并讨论了所用训练方法的稳定性\cite{self-challenging-2506}。这类观察提醒我们，任务能够执行检查不等于训练过程必然稳定。但从曲线先升后降，还不能独自定位到优化器、数据构成、奖励实现或更新步长中的某一项。需要保存训练配置、中间检查点和固定评价，才能进一步分辨原因。本书没有独立复现实验，不把外部报告扩大成“所有约束都无效”的结论。

\end{sloppypar}

\begin{sloppypar}
应把环境变化、更新后的能力变化和评价口径变化分开。对同一旧版本在新环境中重新评价，有助于观察环境因素；把新旧版本放回共同环境，有助于观察版本因素。回滚可以恢复某个已知版本，但不保证环境、外部副作用和优化状态一并恢复。发生在更新算子内部的故障同样可以用检查点阻止继续发布；“能够止损”和“已经解释原因”是两件事。

\end{sloppypar}

\begin{sloppypar}
本书对遗忘的观测沿用两条现成定义：遗忘量 $FGT_t$ 与后向迁移 $BWT_t$（2507.21046 §7.1.2 的定义式\cite{self-evolving-survey-2507}；原论文给的是定义，不是任何方法的数值）。
应用这些定义需要说明各版本在哪些共同任务上被评价，并保存可比较的检查点。回归集可以预先封存，也可以后来独立采集；关键是不能按已见的处理结果挑选有利题目，再将选择后的成绩当作预先约定的保证。新采集集回答的是对应新范围的问题，不能追溯冒充原来的评价。P6的E12尚未实跑，当前没有来自该受控实验的退化与恢复估计。
\end{sloppypar}

\subsection{遗忘与后向迁移的参照不同}

令 $J_{i,t}$ 表示完成第$t$项学习后在旧任务$i$上的得分，数值越大表示表现越好，并要求各检查点沿用可比评价。若使用误差或费用，可先在声明的尺度上取其相反数，再按同一方向解释。跨任务平均还需保持指标与尺度可比。对整数$t\ge2$，所引综述定义
\begin{equation}
 FGT_t=\frac1{t-1}\sum_{i=1}^{t-1}\left(\max_{j\in\{i,\ldots,t\}}J_{i,j}-J_{i,t}\right),\qquad
 BWT_t=\frac1{t-1}\sum_{i=1}^{t-1}(J_{i,t}-J_{i,i}).
 \label{eq:ch16-retention}
\end{equation}
前者比较历史最佳与当前，后者比较初学该任务时与当前，因此二者通常不是同一个量换符号。$BWT_t>0$表示在这些旧任务上的平均后向增益；$FGT_t=0$只表示当前没有低于所记录的历史最佳，不保证所有未测能力不变。该定义来自原文第7.1.2节\cite{self-evolving-survey-2507}，不是本书已有的长期运行估计。

用一个共同评价下的教学例子区分参照。完成第三项学习后，任务一的三个记录依次为\calcnum{0.6}、\calcnum{0.9}、\calcnum{0.8}，任务二从初学时的\calcnum{0.6}变为当前的\calcnum{0.7}。按式~\eqref{eq:ch16-retention}，两项旧任务的遗忘量平均为$[\calcnum{0.1}+\calcnum{0}]/2=\calcnum{0.05}$，后向迁移平均为$[\calcnum{0.2}+\calcnum{0.1}]/2=\calcnum{0.15}$。后向迁移为正与遗忘量为正可以同时出现：当前仍优于初学时，却已低于某次中间检查的最好表现。这些数值均为构造，不是新的持续学习实测。

% ---------- 5. 反馈瓶颈 ----------
\subsection{长期运行：从成绩曲线到状态过程}
\label{sec:ch16-longrun}

\subsubsection{有界、单调和收敛不是同一个条件}

通过率天然位于零与一之间，这只说明数值有界。构造序列
\[
 q_t=0.5+0.1(-1)^t
\]
会在两个值之间往返，虽然每一步都处在合法区间内，却没有单一极限。这是数学教学反例，不是某个Agent的实测曲线。它说明“评价器固定，所以系统收敛”缺少条件。评价器固定只能使数值指向相同测量方式；更新算子仍可能振荡、忘记、随机探索或者反复切换策略。

若记录的是历史最佳
\[
 m_t=\max_{0\leq j\leq t} q_j,
\]
则这个序列按定义不下降。单调不减且有上界的实数序列必有有限极限，但这里的结论只涉及所记录的最佳分数。它没有说明当前部署版本不下降，没有说明最佳候选在独立复测时仍保持同样分数，更没有说明每一步都产生了新能力。第12章展示的最佳曲线与逐候选曲线需要同时解释，原因就在这里。

公开轨迹也可能同时包含初末净收益与中途退化。SEAGym预印本在DeepSeek-V4-Flash上的五轮AHE设置中，对同一组\cited{80}道训练任务进行回放：初态通过\cited{34}题，第四轮后降至\cited{6}题，末态恢复至\cited{43}题（\cite{seagym-2606} §4.1、§4.3）。只看初末差值，就会遗漏这段退化与恢复。这里回放的是已见训练任务，不能把这条轨迹当作封存迁移成绩，也不能据此估计多次整环运行的稳定性或长期故障率。它说明为什么当前版本、历史最佳和中间检查点值得分别保存。

有限前缀也不能唯一决定未来。一段已经观测到的上升曲线，可以与之后的平台期相接，也可以与之后的下降相接；若评价没有变化，这些延续仍可能对应不同的学习过程。要主张某种未来趋势，需要额外模型假设、跨时间的验证或对更新机制的约束。把前几轮的斜率外推到任意长时间，只是在作预测，不是从有限记录直接读取事实。

因此，“到这一轮停止”应是一条明确规则。可以因为预算用完而停，可以因为连续若干轮没有可接受候选而停，也可以因为恢复风险超出发布阈值而停。三者表达不同的终止原因。第一种不能自动写成方法已经饱和，第二种取决于候选生成与筛选能力，第三种说明当前系统不宜继续发布而不一定说明搜索完全没有价值。

\subsubsection{平均收益之外，还要关心暴露时间}

长期系统反复执行更新，风险也会沿时间累积。为说明这一点，构造一个极简教学模型：每次发布在此前未出故障的条件下，发生指定故障的概率恒为\calcnum{0.001}。假设这个条件概率在所讨论序列中不变，则连续$N$次都未出现该故障的概率为
\[
 P(\text{无故障至}N)=\prod_{t=1}^{N}(1-h_t),\qquad h_t=0.001.
\]
这来自条件概率的乘法规则；真实系统的$h_t$可以随状态和历史改变，不能直接套用常数。对上述人为设定，\calcnum{100}次后的无故障概率约为\calcnum{0.9048}，\calcnum{500}次后约为\calcnum{0.6064}。它不是本书测得的可靠性，也不表示所有更新都有相同风险。

这个例子说明，单次风险看似很小，仍须与暴露次数一起解释。更频繁的更新可能更快适应，也可能增加触及失效路径的机会。第13章的触发器因此不能只优化即时成功率；需要同时记录更新次数、候选淘汰、恢复动作和停止原因。如果某个策略只因更早耗尽预算而少发布，其故障次数较少不一定意味着每次发布更可靠。

零次观察到故障也不是零风险。若另构造独立同分布的伯努利检查，观察正整数$n$次均未失败，对固定故障概率$p$的单侧\calcnum{95\%}上限可以通过
\[
 (1-p)^n=0.05
\]
得到$p_{\mathrm{upper}}=1-0.05^{1/n}$。在\calcnum{30}次检查时约为\calcnum{9.50\%}，在\calcnum{300}次时约为\calcnum{0.994\%}。这些数字只用于说明零故障记录的精度，不是本书的发布门槛；相关任务、选择性保留和变化中的系统都可能违反该模型前提。

对于Agent，更实际的记录是条件化的风险描述：在什么版本、什么权限、什么任务和什么恢复能力下，没有观察到哪类故障。一个无外部写入权限的离线执行器，与一个能修改共享状态的部署执行器，即使通过率相同，也不能共享完整的风险结论。新权限应成为新评价条件，而不是隐藏在版本说明里。

\subsubsection{恢复不是把分数改回去}

需要恢复的系统状态至少包含当前执行器、持久经验、搜索档案、未完成任务、资源账和外部环境。恢复执行器代码只回到了其中一部分。若某条坏经验已经被压缩进新摘要，删除原始条目未必移除了它的后续影响；若一次操作已经改变外部记录，恢复本地文件也未必撤销该操作。恢复声明应列明覆盖哪些状态，以及哪些效果只能补偿、不能直接撤销。

发布与恢复之间还可能存在时间差。假设旧版本已做完一批任务，新版本正在处理另一批，发现回归时不仅要选择以后使用谁，还要处理新版本已产生的中间产物。若评价只在最终汇总时进行，那么发现问题的延迟就可能影响需要复核的工作量。低频评价节省成本，但它的代价包括更长的故障暴露时间；这也是第14章把触发、隔离和回滚分开讨论的原因。

一个可审查的停止记录，应同时写明触发信号、被隔离的版本、恢复对象、未恢复部分和后续核验。记录“已回滚”不足以支持“系统已恢复正常”。后者需要回到共同任务或明确的新环境中重新检查。在目前没有真实恢复试验的地方，本书保留设计要求而不填入成功率。

\section{元进化、开放式学习与可证伪的研究问题}
\label{sec:ch16-meta}

\begin{sloppypar}
第12章的DGM与HyperAgents说明，可以允许系统修改部分自身实现，同时保留外部任务和评价参照\cite{dgm-2505,hyperagents-2603}。修改器可变不意味着所有评价都必须由它自行定义；有归档也不等于所有后代都好于父代。要研究更开放的设置，可以扩大允许修改的对象，但每扩大一步，都应说明哪些比较仍沿用共同尺度、哪些问题需另设测量，而非笼统宣称“自指无法评估”。

\end{sloppypar}

\begin{sloppypar}
错误的持续时间由实际写入与传播路径决定。一次推理也可能触发持久的文件写入或外部操作，不会必然随上下文结束而消失；参数可以回滚的前提是保存了可用检查点；共享技能可能扩散的前提是存在传播与加载路径。应记录状态在哪一层被改变、哪些消费者已经读取、能否撤销、撤销要花多大代价。仅按“推理、参数、技能”三个名字排列严重程度，会遗漏真正需要恢复的对象。

\end{sloppypar}

\begin{sloppypar}
权限问题也应落实到具体路径。系统能读一段护栏文字，不等于能改宿主进程的权限；可以修改任务代码，不等于可以覆写独立保存的评价记录。反过来，如果分数解析、日志生成和执行脚本都属于可修改面，冻结题目仍不足以保护测量链。第12章据DGM的完整性案例讨论了这一点\cite{dgm-2505}。本书把代码边界、执行凭据、外部日志和结果复核分别列为检查项，不把所有运行时内控制一概宣判为无效。

\end{sloppypar}

\begin{sloppypar}
最后一件必须说清的是可证伪性。上面引用的综述在 §8 列了六个编号开放问题（按杠杆粗排），
这些是研究议程，不是本书已经执行的检验。本书的分工是把其中的一部分写成带 falsifier 的命题：
接地兑换率写成了候选命题 P8 与实验 E02b（\S\ref{sec:ch14-grounding}），
稳定性工程拆成了 P6 与 E12，可信累积拆成了 P3 / P3b 与 E03 / E03c。
剩下的部分——尤其是“系统能否察觉目标本身已经陈旧”——本书明写为开放问题，不假装有答案。
\end{sloppypar}

% ---------- 7. 四类闭环退化 ----------
\subsection{本书结论的有效范围：bounded 与 open-ended 的边界}
\label{sec:ch16-bounded}

\begin{sloppypar}
这里集中登记本书的整体适用范围，与各章的局部限制一起阅读。
\end{sloppypar}

\begin{itemize}
  \item \textbf{有界自改进（bounded self-improvement）}：本书操作上指在声明的任务、更新对象、资源与评价范围内考察改进。固定外部评价有助于比较，但不意味着序列必然收敛，也不证明所有相关能力已被覆盖。
  \item \textbf{开放式自我改进（open-ended RSI）}：涉及长期扩大任务、能力或研究目标的更强主张。是否修改评审器是一个设计轴，是否保留外部锚点是另一个轴；两者不能被压成互斥的定义。

\end{itemize}

\begin{sloppypar}
本书与上述综述使用这些词时，需要区分文献分类与实际检验。综述将其语料的主要工作归为有界改进，这是作者在该分类下的判断。本书自己的限制更具体：当前模型试验使用有限的程序任务、限定调用预算和指定验收规则，没有测试长期目标演化，也没有估计无限时间下的能力趋势。开放式系统仍可接受有限阶段评价，但有限阶段结果不能自动支撑无界增长主张。

\end{sloppypar}

\begin{figure}[htbp]
\centering
% Operational evaluation scopes; not a convergence theorem.
\begin{tikzpicture}[
 box/.style={draw,rounded corners=2pt,text width=3.6cm,minimum height=1.25cm,align=center,font=\small},
 arr/.style={-{Stealth[length=1.8mm]},thick}]
 \node[box,fill=blue!8] (a) at (0,0) {本书已观测范围\\有限任务、预算与更新对象};
 \node[box,fill=green!8] (b) at (4.7,0) {扩大可修改范围\\保留或重建共同评价};
 \node[box,fill=orange!10] (c) at (9.4,0) {长期与目标演化\\另定义范围、状态与资源};
 \draw[arr] (a) -- (b);
 \draw[arr] (b) -- (c);
 \node[align=center,text width=3.8cm,font=\footnotesize] at (0,-1.65)
 {固定规则有助比较\\仍须审查实现与覆盖};
 \node[align=center,text width=3.8cm,font=\footnotesize] at (4.7,-1.65)
 {修改器或评审器可变\\不等于无法外部评价};
 \node[align=center,text width=3.8cm,font=\footnotesize] at (9.4,-1.65)
 {有限阶段可以测量\\不能由此证明无界增长};
 \draw[dashed,gray] (-1.9,-2.55) -- (11.3,-2.55);
 \node[align=center,text width=13cm,font=\footnotesize] at (4.7,-3.15)
 {每次扩展都需要新证据。固定评价、数值有界、历史最佳单调与系统收敛是不同条件。};
\end{tikzpicture}

\caption{从本书有限比较到更强主张所需的评价扩展。这是证据范围图，不是系统必然经历的升级阶梯，也不以是否修改验证器独自定义开放性。右侧范围尚无本书实测。}
\label{fig:ch16-scope}
\end{figure}


\begin{sloppypar}
对于长期主张，至少要交代三个问题：后续有效信息从何而来，已有经验如何避免被错误替代，以及资源如何响应任务难度增长。它们是本书的审查问题，不是一条已经证明每个系统都受同一数值约束的定理。理论结果应保留各自假设，经济估计应保留样本和模型设定，有限实验应保留任务范围。把三类证据拼接起来，可以形成研究计划，不能因此得出它们共同证明了某个未来轨迹。

\end{sloppypar}

\begin{principlebox}
\textbf{边界声明}：本书的实测主张限于已声明的有界设置。
具体地：本书模型试验以指定的 $\mathcal V$ 为比较参照，设计上不将任务划分与验收规则列为更新对象；但实现审计发现历史 E01/E02/E03 允许读取受保护测试文件，部分经验生成载荷也包含测试材料。因此，禁止修改与禁止读取必须分开检查，不能把设计意图当成已实现的封存边界。
若允许系统修改自己的验收器，旧运行数字只能作为原设置的记录，不能直接外推。仍可另外设置独立外部评价，但必须重新规定收益的参照和数据边界。
第14章的污染注入实验位于这条线的内侧：噪声臂实际翻转了\measured{2}条反馈标签前缀，原始测试输出仍然保留。它相对无记忆基线的通过率差为\measured{$-0.1500$}；相对保留原标签的测试反馈臂，差为\measured{$-0.0500$}，95\%题级自助区间为\measured{$[-0.2000,+0.1125]$}。两个反馈臂分别完成训练与经验构建，因此后一差值比较的是本轮两套流程的结果，尚未单独识别标签翻转的作用。这次处理也没有改写验收器；区间同时兼容正负方向，不能据此断言污染普遍有害或无害。
\end{principlebox}

\begin{sloppypar}
这条声明对阅读本书的方式有一个直接后果：\textbf{不得把本书的负结果读成“自进化无用”}。
本书用已有案例说明局部机制通路，并以问题表明确后续可以研究的条件。同理，本书的边界声明也不是在说开放式一侧不可能，
而是在说它超出了本书的测量范围。
\end{sloppypar}

% ---------- 9. 方向生成 ----------
\section{闭环退化的四种形式}
\label{sec:ch16-collapse}

\begin{sloppypar}
第\ref{ch14}章\S\ref{sec:ch14-collapse} 讨论了前三类退化，并另述框架锁定，本章补的是它们的分界线：
不同退化需要不同的观测量，任何单个签名都不保证充分检测。即使监测器发现变化，也未必已经区分原因；看不到变化则可能因为没有记录相关维度。
\end{sloppypar}

\begin{table}[htbp]
\centering
\caption{四类闭环风险：候选信号与检测局限}
\label{tab:ch16-collapse}
\begin{footnotesize}
\begin{tabular}{>{\raggedright\arraybackslash}p{2.2cm}>{\raggedright\arraybackslash}p{3.4cm}>{\raggedright\arraybackslash}p{3.0cm}>{\raggedright\arraybackslash}p{2.6cm}>{\raggedright\arraybackslash}p{2.2cm}}
\toprule
退化形式 & 机制 & 环内签名 & 干预点 & 检测局限 \\
\midrule
自证循环 & 输出与评审的相关偏差被选择反复保留 & 信心上升而独立正确率不升 & 校准评审；引入异源核查 & 共享权重不是必要条件 \\
\midrule
输出分布退化 & 部分模式或尾部信息在更新中丢失 & 覆盖变化与分组误差 & 保存参考数据；检查采样与训练 & 熵不足以唯一诊断 \\
\midrule
课程收窄 & 候选选择集中到少数题族 & 族占比、覆盖与固定集成绩 & 外部题源；固定覆盖检查 & 先升后降不是收窄的专属信号 \\
\midrule
框架锁定 & 原目标与当前需求逐渐脱节 & 需求、结果与原分数之间的差异 & 引入需求证据；复核目标 & 只记录旧分数可能看不见 \\
\bottomrule
\end{tabular}
\end{footnotesize}
\end{table}

\begin{sloppypar}
框架锁定（frame lock-in）值得单独说一句，因为它最能说明“验证器冻结”这条公理的两面性。
一个可以接地、诚实、多样且稳定的闭环，仍然可能在优化一个已经过期的目标（2607.07663 §5.3 失效模式 4）\cite{rsi-survey-2607}。
这个风险与前三类并不互斥。若系统只跟踪旧目标的分数，就不能从该分数单独推断目标仍合适；若另外记录需求、用户后续结果或业务损失，便可能观察到脱节。能否检测取决于可用信息和设计，不能作环内永远无信号的绝对断言。
本书第 13 章讨论“何时更新”，框架锁定问的是“更新所服务的目标是否已经过期”——
\textbf{两件事不能合并}，前者的答案再好也回答不了后者。
\end{sloppypar}

% ---------- 8. bounded vs open-ended ----------
\section{方向生成：验证层级索引不到的那一半}
\label{sec:ch16-direction}

\begin{sloppypar}
第4章区分形式验证、执行反馈、学习式评审和内在信号。所引综述的层级是定性观察，不是任何任务上的可靠性全排序。\textbf{方向评估}可以借助这些信号检查候选的某些性质；是否值得投入还涉及目标与资源，不能仅凭信号类型决定。综述举出了针对研究提案部分质量的基准尝试，而非宣称所有方向价值都已有完备测量。
\end{sloppypar}

\begin{sloppypar}
另一半是\textbf{方向生成}：哪些问题值得被评估，以及一个仍然被高分评价的问题什么时候已经不值得回答。
它在层级\textbf{之前}，不是层级的一级（2607.07663 §5.4）\cite{rsi-survey-2607}。
一个闭环可以爬到层级顶端、保持诚实的评估器、检测出自己的代理指标被腐蚀，
却仍然在一个早该放弃的问题上做得很好——\textbf{因为层级把“问题”当作给定的}。
\end{sloppypar}

\begin{sloppypar}
这件事对“人在环上”的长期含义是具体的，不是一句“人类仍需参与”。
人在这一侧的工作不是给更细的分数，而是提供一个不位于闭环判据之内的视角；
新的视角也可以通过新增观测、工具、交互或候选目标进入系统；它不是人类独有的形式保证。关键是新证据有没有改变问题设定，而不是参与者被叫作人还是模型。
\end{sloppypar}

\begin{sloppypar}
综述没有给出统一的测量标准，而是提出了候选可观测量（2607.07663 §8-(6)）\cite{rsi-survey-2607}：
\textbf{（一）}productive abandonment——放弃一条局部指标仍在改善的路线，以后来的验证打分；
\textbf{（二）}invariant transfer——系统在重构自身任务时，把哪些约束带过了重构；
\textbf{（三）}absence-sensitivity——一个缺失的结果是否被当作证据；
\textbf{（四）}reframing latency——环境漂移之后，一个闭环需要多少证据才放弃一个陈旧目标。
这些是该综述提出的候选测量方向。本书未对截至交付日的所有基准作穷尽检索，因此不沿用“当前不存在任何这种基准”的全称断言。
\end{sloppypar}

\openquestion{本书对方向生成没有任何实验，也不打算在首版里补。
原因不是它不重要，而是它需要先有一个“目标已经陈旧”的可判定定义，
而本书尚未冻结这种任务设计和评分规则。第 13 章的触发信号、第 14 章的漂移检测都只在给定目标内部工作；
把它们当成“系统会察觉目标过期”的证据，是本章要挡住的最后一个错误。}

% ---------- 10. 判决 ----------
\subsection{改变目标时，怎样保留可比较的历史}
\label{sec:ch16-goalchange}

\subsubsection{一个教学案例：关闭工单的速度越来越快}

设想一个服务Agent，最初的目标是在规定时间内完成工单，验收器统计关闭比例。系统学会模板回答后，关闭速度持续提高。但需求后来变化，越来越多工单需要跨系统核对，用户又会在被关闭后重新提交。同一个关闭指标仍然上升，实际问题却没有得到同样程度的解决。这个场景是本书构造的教学例，不是所引用论文中的真实产品测量。

诊断应先排除几个不同问题。若系统在没有完成必要操作时标记完成，这是执行与验收脱节；若原本有重开记录但汇总器漏算，这是测量实现错误；若规则从一开始就只要求关闭，而业务现在更在意问题最终解决，则涉及目标变化。三个问题可能同时存在，却需要不同修复。只给出“奖励被破解”这个标签，会掩盖应该改规则、改代码还是改需求定义。

若增加“后续重开率”，系统就多获得一个观察窗口，但不自动得到新的总目标。把关闭率与重开率加权，需要选择时间跨度和权重；较晚重开的工单还存在观察不足。更进一步，减少重开也可能被错误地优化为阻止用户重新提交。因此新指标应对应可追踪的结果，而不仅是容易计数的新动作。目标设计本身也需要反例，不是旧指标失败后随手增加一列即可完成。

这说明框架锁定未必完全没有环内信号。只要系统原来就能观察用户后续行为，它可能发现关闭率与解决结果脱节；若信息从未进入系统，仅靠旧分数则难以识别。真正的问题是哪些信息可见、系统是否会把冲突保留为问题，以及谁有权限修订发布规则。把人加入流程也不保证发现问题，若人同样只看到关闭率仪表盘，信息缺口仍然存在。

\subsubsection{目标版本和系统版本要分开编号}

将旧目标的评分规则记作$\mathcal O_0$，新目标的评分规则记作$\mathcal O_1$，旧执行器记作$A_0$，新执行器记作$A_1$。这里用$\mathcal O$区别于表示工作流结构的$G$。若直接比较$\mathcal O_0(A_0)$和$\mathcal O_1(A_1)$，差异同时包含系统变化与尺度变化。更有解释力的做法是在能够重放的共同任务上保留四格：两个执行器分别按两个目标评价。这样可以看见，新版本是否仅适应了新口径，以及旧版本在新口径下原来是什么状态。

四格并不总能补齐。例如，新目标需要记录过去从未保存的长期结果，就无法从旧日志重建。此时应把历史不可比的部分留下，而不是凭旧输出猜测新目标得分。可以从共同起点重新运行新旧版本，或者从新目标启用日起建立新的时间序列；无论选哪种，都要说明代价和解释范围。不能用一条连续折线把不可比的前后阶段连接起来，再称其为持续进步。

目标修改也不意味着所有约束都需要修改。工单场景中，解决目标可以变化，但访问权限、不可删除的记录、用户明确表达的限制可能仍需保留。应该分别登记“优化什么”和“不能破坏什么”。若新目标只有通过放松约束才能取得高分，应将这种变化作为处理的一部分报告。否则看似能力上升，实际可能只是验收变宽。

评价器更新后，原有记忆和技能也需要重新解释。一条“先关闭再补信息”的旧经验可能在旧规则下获准，却与新目标冲突。仅替换评分脚本而不检查持久物，系统仍会沿旧经验行动。迁移检查因此应覆盖旧经验的适用条件，而不只是模型版本。能够删除不再适用的经验，也可能是目标转移中的有效行为，不应把记忆条数减少一律视为能力退化。

\subsubsection{提出方向与批准方向之间需要留出检查}

系统可以生成候选目标，可以给出需求变化的证据，也可以模拟某项规则修改的后果。这些能力都值得研究，但“提出建议”与“直接替换正式目标”应分开记录。前者允许探索多个解释，后者会改变所有后续成绩的意义。如果一项候选规则因提高本系统排名而被自动采纳，目标选择与结果选择就重新耦合了。

本书对方向生成的可测问题，可以从一个窄范围开始：给定明确的环境变化记录，系统能否识别旧目标漏掉的结果，提出有证据支持的修改，并保留指定的不变约束。这里并不声称覆盖科学品味或所有价值判断；它只是把一个宽泛愿望缩成可检查任务。需要分别评价察觉、解释、提案与后果，不能只凭它说出“目标需要调整”就认定方向能力成立。

在这种设计中，拒绝修改也可能是正确结果。环境出现短期噪声，不一定意味着目标陈旧；为了追求新颖而频繁改变目标，会让历史证据失去作用。应设置无需修改的对照情形，记录误报与漏报，避免只考察那些必然需要改写的问题。只让系统处理“旧目标一定错”的任务，会奖励一种不加分辨的改写倾向。

\section{十个研究问题：已有观察与下一步设计}
\label{sec:ch16-verdicts}

研究议程从可以观察的机制出发。表\ref{tab:ch16-verdicts}列出十个问题：现有材料在哪些条件下提供了什么信息，下一项对照又能回答什么。P编号沿用历史研究索引，便于读者回到协议与分析文件；方向性的原假设及旧状态保存在附录\ref{sec:appg-history}。本表不把新问题追溯声称为旧试验的预注册目标。

\begin{small}
\setlength{\tabcolsep}{4pt}
\begin{longtable}{>{\raggedright\arraybackslash}p{0.65cm}>{\raggedright\arraybackslash}p{3.0cm}>{\raggedright\arraybackslash}p{5.0cm}>{\raggedright\arraybackslash}p{3.4cm}}
\caption{研究问题、已有观察及下一项可识别的比较}\label{tab:ch16-verdicts}\\
\toprule
ID & 研究问题 & 已有观察及适用范围 & 下一步设计 \\
\midrule
\endfirsthead
\multicolumn{4}{l}{表\thetable（续）}\\
\toprule
ID & 研究问题 & 已有观察及适用范围 & 下一步设计 \\
\midrule
\endhead
\midrule\multicolumn{4}{r}{续下页}\\\endfoot
\bottomrule\endlastfoot
P1 & 在什么任务与预算条件下，持久更新优于任务内重试？ & E01：格4减格1为 \measured{+0.0500}，区间 \measured{$[-0.1000,+0.1750]$}；格1未实现重试后选择，完整总成本未对齐 & 匹配重试、选择与完整预算 \\
P1b & 共享同一起始记忆后，继续更新比只读复用多带来什么？ & E01：格4减格2为 \measured{+0.0625}，区间 \measured{$[-0.1375,+0.2750]$}；两臂经验来源不同 & 同源记忆与只读/更新对照 \\
P2 & 怎样区分反馈质量与更新对象各自的贡献？ & 统一交叉设计：旧E01/E02循环不同，标准化不能排除混淆；需直接操纵反馈和更新对象 & 共同循环内交叉操纵两因素 \\
P2b & 不同反馈机制在什么条件下提供有效改进信号？ & E02：自评减执行为 \measured{$-0.0875$}，区间 \measured{$[-0.2625,+0.0875]$}；单元测试臂不能代表形式证明 & 真实对应的反馈操作与联合比较 \\
P3 & 记忆形式与容量如何共同影响跨任务迁移？ & E03：不含冲突基线的直接比较，结构化减摘要为 \measured{$-0.0125$}，区间 \measured{$[-0.1500,+0.1125]$}；缺容量对照且仅两个留出族 & 同源材料上的形式与容量交叉设计 \\
P3b & 拟合目标怎样影响经验跨模型与跨执行接口的迁移？ & E03c领域拟合减求解器拟合：同模型差约\measured{0.0000}，区间\measured{$[-0.1375,+0.1500]$}；跨模型差\measured{$+0.0625$}，区间\measured{$[-0.0625,+0.1750]$}；两侧均跨0，且只换过一个替代模型 & 读取者各自基线、同期版本与接口变化 \\
P4 & 在给定质量和成本条件下，技能复用何时回本？ & E05：v2完成单族流程试跑，发现格式失败；未获得有效技能收益与成本交点证据 & 有效技能、复用次数与完整成本 \\
P5 & 不同任务流中，怎样选择更新时机？ & E11：未实跑；已复现信号时点与重复切段错误，v2时机、分析与总预算运行器已有实现和离线检查记录，尚无新模型轨迹 & 共同任务流与总资源下比较三策略 \\
P6 & 哪些更新会引起旧能力退化，哪些条件有助于保持与恢复？ & E12：未实跑；v2已修订固定评价、候选池与恢复分母，装置检查已有记录，尚无能力与恢复实测 & 固定评价下区分内容、幅度与恢复 \\
P7 & 元层改进能迁移到哪些新任务，其贡献如何识别？ & E10：本书E10未跑；HyperAgents v2已有任务/元组件交叉对照，在其数学评分配置中主要增益随元Agent迁移 & 扩展目标任务族，继续分离元层内部机制与资源作用 \\
\end{longtable}
\end{small}

表中读数来自E01的\texttt{20260923T152230Z-90acda}、E02的\texttt{20260924T153511Z-421eac}和E03的\texttt{20260925T152453Z-254845}运行，沿用历史题级自助区间。它们描述所测任务与实现，不自动推广到任意新任务族。P2b关联的反馈层级在原综述中属于定性观察，因而这里按反馈机制与条件提出问题。

E11、E12的离线装置与协议为研究更新时机和恢复提供执行基础；策略收益与行为恢复仍须在真实任务上比较。工程检查、模型观察和因果解释各有用途，分别进入相应证据链，实验边界见附录\ref{app:lab}。

\subsection{更多题不能修复比较对象不匹配}

P3问形式与容量如何共同影响迁移，当前E03只改变形式，没有交叉操纵容量。
因此即使把每个现有臂跑得更精确，也只能更精确地回答已实现的形式比较。
要比较两种因素的影响，必须先规定它们的可比较变化范围，再安排相应对照。
交叉设计改变能识别的作用；增加原有题目则主要改变已测比较的精度。

P2 的问题也不能通过两个标准化效应量相除解决。
不同执行循环、反馈权限和对照对象共同变化时，标准化只改变数值单位，不产生因果可比性。
下一轮应在共同任务、执行器和预算下安排反馈与更新对象的交叉设计。
如果分母接近零，效应量之比还会不稳定，故主分析量本身也须在运行前确定。

历史P6使用“任何更新”的普遍量词，原句见附录\ref{sec:appg-history}。新的问题先把更新范围、保持指标和恢复对象规定清楚。
对特定更新范围，可以报告退化分布、回归维度和恢复成功率；
若要提供反例，则要有可信地排除相应退化的比较，而不能将“不显著下降”当作“没有下降”。
本书统一下降量为旧版本得分减新版本得分，正值才表示退化，避免旧表中的方向混淆。

\subsection{历史样本量表实际上规划了区间精度}

旧稿把表~\ref{tab:ch16-power} 称为判决力表，使用的公式却没有目标功效。
现保留计算值，将它们正确标注为独立题假设下、目标半宽为 $0.03$ 的精度近似。
这些值描述后续规模的一个参考，不能直接成为确认性检验的充分条件。

\begin{table}[htbp]
\centering
\caption{试验方差对应的历史精度估算，不是正式功效设计}
\label{tab:ch16-power}
\begin{footnotesize}
\begin{tabular}{p{4.2cm}p{2.1cm}p{2.2cm}p{2.3cm}}
\toprule
比较 & 配对差标准差 & 精度估算题数 & 现有题数 \\
\midrule
E01 格4减格1 & \measured{0.2875} & \calcnum{353} & \measured{16} \\
E01 格4减格2 & \measured{0.4425} & \calcnum{836} & \measured{16} \\
E02 自评减执行 & \measured{0.3722} & \calcnum{592} & \measured{16} \\
E02 测试验证减自评 & \measured{0.3183} & \calcnum{433} & \measured{16} \\
E03 结构化减摘要 & \measured{0.2778} & \calcnum{330} & \measured{16} \\
\bottomrule
\end{tabular}
\end{footnotesize}
\end{table}

\texttt{calculations/audit\_sample\_size.py} 用未舍入标准差复算全部历史值，
并另外给出双侧显著性水平 $0.05$、功效 $0.80$、检出差异 $0.03$ 的规划示例。
在其独立题与正态近似假设下，E01 两个比较分别约需 \calcnum{721} 与 \calcnum{1708} 道题。
这些数不含族相关、多重比较和方差估计不确定性的调整，也未被追溯宣称为原运行的预注册规模。
方法解释见第 3 章和附录~\ref{app:stats}。

\subsection{补齐一条实验臂之后还要检查什么}

E03c 的缺失臂已于 2026-09-28 补齐，两条拟合目标在新读取者上的试验比较现在都有了。
补齐带来的三样东西是：跨模型的配对差（点估计 \measured{$+0.0625$}、区间 \measured{$[-0.0625,+0.1750]$}）、
领域拟合一侧逐题点估计反号的题数（\measured{2}/\measured{16}）、以及换读取者的成本对照。
回答P3b还需要读取者范围与同期比较：
区间跨0，且本书只换过\emph{一个}替代模型。历史\calcnum{16.7}倍是区间半宽近似规划的比例，不是已经知道的确认性样本缺口。
跨日期续跑需要保存父运行、记忆哈希、配置和时间信息；相同服务别名也不保证模型修订未变。
如果本轮仍只覆盖少量固定任务，就应保留小样本试验的定位，并为确认性比较设置同期基线。

研究问题表用于选择下一条有信息价值的证据：补臂消除矩阵缺项，
新对照改善归属识别，更多独立族扩大覆盖，合适的规模和分析提高精度。
把这些工作分开，读者才能判断一笔新增实验预算究竟买到了什么信息。

扩展问题P8依赖尚未执行的E02b，单独作为反馈接地频率的研究方向保留。
第 14 章保留其研究问题，本章只登记它尚未成为本书已执行的检验。


\subsection{研究预算应该先消除哪一种不确定性}
\label{sec:ch16-information}

\subsubsection{识别、精度与覆盖有不同的购买方式}

识别缺口意味着观察差异还不能归到所问处理上。此时扩大同一种混合比较，通常只会得到更稳定的混合结果。P2需要共同循环内的反馈与更新对象操纵，P3需要形式与容量的交叉变化，不能用更多原有题目替代这些对照。新增对照的价值是改变可以回答的问题，而非仅使误差条更短。

精度缺口则是在比较对象已经明确后，数据还不能把重要效应范围区分开。新增重复可以减少一部分随机性，但如果所有重复都集中在同一题族，它不必然增加对新族的精度。应把族、题与运行层级分开，说明预算增加在哪一层。本书历史题级区间不能被重新命名为跨任意任务族的区间；要回答后一种问题，需要相应抽样和分析。

覆盖缺口关心结论适用于哪些条件。多增加一种读取者、一次环境变化或一类失败模式，可能比同一条件下重复更多次提供不同的信息。但少量新增条件也不构成全面覆盖。选择哪种条件，应跟随主张最可能失效的边界：技能是否依赖特定接口，记忆是否依赖特定读取者，修改器是否依赖源任务的评价结构。选择困难反例有助于探索边界，但不能把这种有目的的样本当作总体随机样本。

装置缺口又不同。若成本重复计数、候选未真正准入或统计切段错误，继续付费运行可能只会增加不可解释的数据。离线检查先修复这些问题，是减少浪费；它也应有停止点。装置通过既定检查后，不应无限增加镜像测试来替代真实验证。接下来应该运行适当规模的流程试验，观察真实后端是否触及离线脚本没有覆盖的路径。

\subsubsection{确认性分析与发现性分析可以共存}

一次研究可以有预先指定的主要比较，也可以在轨迹中发现新的失败类型。前者按原规则报告，后者清楚标为事后探索。它们不必互相排斥。例如，E05工程试跑暴露格式解析失败，即使原目标是检查成本生命周期，这个新发现仍值得记录。问题在于不能发现一种有利切分后，把它追溯写成从一开始就指定的主要问题。

筛选候选本身会改变看到的结果。构造一个纯教学情形：若有\calcnum{20}个相互独立、实际上均无效的比较，每个都以\calcnum{0.05}概率错误报出阳性，那么至少出现一个阳性的概率为
\[
 1-(1-0.05)^{20}\approx\calcnum{0.6415}.
\]
这个数依赖独立性和所设错误概率，不能套给本书十个相关研究问题所涉及的比较，更不是它们的实际误判概率。它只说明从许多读数中选择最好的一个，再按单次比较的口径解释，会遗漏选择过程。

因此，报告最佳候选时应说明候选数、选择依据和独立复测。对探索中形成的新假设，可以在后续新数据上检验；若只能使用旧数据，则明确它是探索性证据。预注册不是禁止学习，而是保留“哪些问题在看结果之前确定”的时间信息。研究者根据新故障修订协议可以是正确行为，前提是新旧版本不被拼成一个从未实际执行的设计。

\subsubsection{命题中的量词决定怎样的反例才够用}

“某些更新能够改善表现”“平均更新能够改善表现”与“每次更新都改善表现”有不同的逻辑结构。一个可信的成功案例可以为存在性提供局部证据，却不足以支持平均主张；一个可信的失败案例可以反对相应范围内的全称保证，却不意味着平均效应为负。命题的措辞应在看到结果之前就确定，不能成功时使用存在量词、失败时又悄悄改成平均量词。

若把P6的“旧能力退化”进一步拆成多个维度，就需要预先说明维度集合。如果把未观测、未来才提出的任意维度都包括进来，有限实验很难形成有用反例；若把范围限定为预先列出的任务质量、费用和恢复指标，就可以研究这组指标上的变化。但这种可操作化缩小了命题，必须公开登记，不能声称新实验已经回答无限范围的原句。

严格相等也不适合用普通不显著性来代替。某维度的新旧差值区间跨零，只说明当前数据与零差兼容；区间仍可能同时包含有实际意义的下降。若工程问题是排除超过阈值的退化，应事先声明允许幅度，并用相应区间比较。阈值应来自应用要求或明确的研究目标，不能在看到区间端点之后调整到刚好通过。

这里还存在成本维度的选择问题。若把任何额外计算都定义为退化，那么一次需要运行更新器的处理几乎按定义就付出成本，原命题容易变成没有判别力的陈述。若研究真正关注质量与总资源的权衡，就应在共同资源约束下比较质量，或把质量约束下的总成本单独列为问题。把定义性代价与经验性副作用分开，才能让实验可能改变我们的判断。

一个实用的命题卡应保留原句、操作化版本、允许反例和未覆盖范围。后续研究若发现原句过宽，可以提出修订，但必须保留旧版本和修订原因。修订命题不是篡改数据；把新命题假装成旧实验最初要检验的命题，才会破坏证据关系。本书将原假设单独归档，并在问题表列出下一项对照；原数据、原目标与新设计因此各有明确身份。

最后，应避免把开放状态当作永远免于反驳的保护层。新证据若在合适范围内形成反例，就应更新判定；若支持了预定效应，也应明确承认支持。公开条件的目的在于让正反两种结论都可能出现，而不是不断移动门槛以保留最初的叙事。

\subsubsection{按信息价值安排下一项比较}

研究问题需要的信息并不处在同一层。一个只有局部点估计的比较、一个尚未隔离因素的比较、一个尚未运行的实验，不能按“看起来接近支持”的程度排成统一排行榜。这种排序会把尚未测量伪装成效果较差，也会把宽区间中的正点估计误当作更有希望。下一步优先级应根据具体可以消除的缺口与成本决定，而非按预期结果给研究问题排主观胜率。

对当前工程链，先修复账目、候选和恢复边界具有明确价值，因为这些缺陷影响随后每次运行的解释。修复完成后，小规模流程试验可以检查真实载荷和失败分母，再决定是否值得扩大。若目标是判定宽泛命题，就不能把流程试验得到的少量回合直接升级为确认性研究；应重新规定任务覆盖、预算配对、统计对象和停止规则。

负向证据同样有优先价值。一个可靠的反例可以推翻相应范围内的全称命题，但不能任意扩大到所有更新机制；一个非劣性比较可以排除超过指定幅度的下降，却不意味着各维度完全不变。结论的力度取决于问题的量词与比较设计，不取决于作者更希望看到成功还是失败。把这条纪律保持到最后一章，才使前面所有机制分析保留可修订性。


% ---------- 11. 条件变化 ----------
\section{把有限结果写成有限主张}
\label{sec:ch16-finiteclaims}

\subsection{四种强度的句子，需要四种不同的证据}

“系统会改进自己”至少可以指四件不同的事。第一件是实现能力：程序确实能把一个候选变化写入持久状态，而且后续执行实际加载了它。第二件是局部效果：在指定任务、预算和评价规则下，采用更新的条件优于一个明确基线。第三件是范围迁移：把版本移到未参与选择的任务族、读取者或环境后，仍保留某种预先定义的收益。第四件是持续性：在更长序列里，更新收益、资源增长和恢复风险形成可接受的关系。它们可以依次成为研究目标，但并不是同一句话的四种修辞。

第15章的静态技能轨迹能支持第一种句子的部分内容：候选技能通过准入，来源代码随后被注入提示。它还不能独自支持第二种，因为后续做对的题也可能在不使用技能时做对；更不能直接支持第三种，因为题族没有改变。动态技能的来源失败则说明，连接“生成候选”和“获得可用持久物”的链条会中断。把成功回合与失败回合一同保存，才能知道目前实现到了哪一层。只展示一段成功的最终代码，会把装置存在、实际采用和效果成立压成一个无法拆开的故事。

局部效果进一步要求回答“相对谁”。相对一次直接回答、相对同样预算内的重试、相对同样信息量的上下文、相对另一种持久表示，是四个不同基线。某个更新方式可以胜过第一种而不胜过第二种，这不构成逻辑矛盾。它可能更有效地利用了额外计算，但还没有证明持久性本身具有额外价值。比较对象应跟随问题确定，不能在结果出来后挑一个最容易超过的基线，再用同一个标题对外宣称普遍有效。

范围迁移还需要说明迁移前保留了什么。如果带走的是整个系统，结果描述整个系统的迁移；如果只带走一条技能，就要固定或重建其余依赖；如果在新环境先重新训练，再报告成绩，就应把适应数据和计算也算进迁移过程。完全禁止适应与允许少量适应都有合理用途，但它们衡量的能力不同。第12章提出的共享元实现对照，就是为了将“整个实现携带了有用知识”与“修改器本身可迁移”分开。

持续性也不是把同一道题多做几次。一个系统可以在稳定任务上逐渐熟练，同时面对新需求时越来越脆弱；也可能短期总分下降，却因探索新任务而积累后续可用的覆盖。必须先规定关注的是固定集合的质量、不断到来的任务流收益、知识覆盖，还是单位资源的改进速度。本书的有限试验没有观察到所有这些量，因此不使用“持续变强”作为它们的共同替代语。

\subsection{用一个结果句保留证据的身份}

一条可以核查的结果句应包含处理、参照、对象、资源、指标与不确定性。它不必写成冗长模板，但缺少其中任何一项，都可能改变读者理解。例如，“这次运行中，带经验条件在指定留出题上的平均通过率高于不带经验条件”只给出了方向。继续补上差值、对应区间和对照限制，读者才知道它是描述性的正向点估计，还是已经得到适当支持的平均收益。

还应区分“存在过一个好版本”和“指定方法通常产生好版本”。前一句的统计单位可以是一个被完整记录的候选，后一句则需要覆盖方法的重复运行、候选选择过程与失败次数。自动搜索保留历史最佳时，最终候选可能表现很好，但挑选它的总预算、未被选中的分支和重新运行的变异性仍是方法的一部分。只把最终版本的推理费用与初始版本比较，会遗漏获得该版本所支付的研究费用。

进一步区分时，还要问重复了哪一个过程。对同一末版反复评价，测到的是该版本在评价条件下的变异；从共同初态重新运行完整更新循环，才开始观察方法生成结果的变异。若主张后期修改器比早期修改器更有效，可以让两者从同一产物、按相同预算生成后继，再用共同规则评价。起始分数相同并不保证起始状态相同。《Auditing the Product, Not the Recursion》预印本的表2、§4.3和表9对这些区别作了系统说明\cite{auditing-product-2610}。整环复跑是否足以支持可靠收益，还取决于重复次数、估计对象与精度；一条上升曲线不能代替这些比较。

“在当前评价器下获得高分”与“满足实际需求”也不能互换。前者可以由运行记录直接核对；后者还取决于指标的含义和覆盖。例如，一段代码通过现有测试说明它满足这些测试的要求，不能由此断言它满足所有未写明的边界情况。这个限制不使测试失去价值，而是要求结果句保留测试范围。把有限检查说清楚，比把它包装成无条件保证更有助于后续复用。

日期和版本是结果句的索引，而非行政附件。模型服务、任务生成器、依赖库、提示模板和记忆内容都可能在两次运行之间变化。只保留一个实验简称，后来就无法判断两个读数属于同一处理。第3章的运行卡与第15章的案例包分别解决汇总和追踪问题：前者说明本次比较是什么，后者让读者能从结论返回具体载荷、产物与验收记录。

\subsection{没有看见的维度，要留为空白}

本书某些表格保留“未运行”“未核查”“费用未知”。这些字段不应被填成零。未记录美元费用不意味着免费，没有保存恢复轨迹不意味着无需恢复，没有遇到某种故障也不意味着故障概率为零。汇总器若把缺失值默认为零，会在成本图上制造一种虚假的优势；若把所有未完成回合删除，又会将装置可靠性从成功率分母中消失。

空白也有不同来源。计划本来没有测量某个维度，属于范围限制；预定要记录却丢失了记录，属于数据完整性问题；由于早期失败而无法进入后续步骤，属于实际过程结果。三者的补救不同。第一种可以设计新实验，第二种可能只能缩小旧结论，第三种应保留失败并分析其后果。不能一律用“后续补跑”把它们混成可以无损续接的空格。

对尚未运行的E10、E11、E12，协议与离线检查可以回答装置是否具有某些预定行为，却不能回答模型处理的真实效果。脚本后端能够精确安排响应，从而检验调用上限、候选准入、切段和恢复分母；它不会提供真实模型的错误分布。证据表应把“装置检查通过”和“效果待测”同时留下。后一格空白不会抹掉前一格工作的价值，前一格完成也不能代替后一格。

\section{条件变化：这些边界什么时候移动}
\label{sec:ch16-conditions}

\begin{itemize}
  \item \textbf{修复对照并扩大独立样本时}：P1/P1b 才能获得与原问题匹配的新证据。
        规模须按新设计计算，达到旧精度公式给出的整数不保证可以判定。
  \item \textbf{补入真正形式证明等处理并重建同源对照时}：P2b才可能获得覆盖原文层级的新证据。
        更换验收器还会改变任务与成本，须同时核对其他层的操作和判据；仅改工具名称或增加样本量都不足以补齐比较。
  \item \textbf{放宽公理 1、允许系统修改 $\mathcal V$ 时}：旧数字不能直接支撑新的更新范围，
        应重建共同评价与版本对应，不能直接复用旧比较（\S\ref{sec:ch16-bounded}）。
  \item \textbf{任务分布由系统自己产生时}：需要分别检查多样性塌缩与框架锁定，
        任务分布熵要从实际接纳记录重算，不能只相信生成器给自己的多样性标签；它仍是可用的局部统计量。
  \item \textbf{经验开始跨模型、跨 harness 复用时}：P3b 必须先报告换读取者后的静态基线，
        否则“经验失效”与“模型更弱”不可分（\S\ref{sec:ch16-eval}）。
\end{itemize}

% ---------- 12. 能说与不能说 ----------
\label{sec:ch16-claims}

本书建立了一套分析自进化Agent的框架：识别被更新的持久组件，追踪反馈与经验的来源，在共同任务和完整资源范围内比较结果，并检查更新后的迁移与恢复。既有试验展示了这套框架如何用于解释局部收益、发现比较失配，以及追查汇总与执行记录之间的冲突。

这些案例形成了明确的机制认识和可执行的研究设计。持久更新何时优于重试、记忆形式与容量如何交互、技能何时回本，以及更新如何影响旧能力，需要在各自明确的条件下继续检验。研究问题表把这些条件连接到下一项比较，使机制分析能够转化为具体的实验选择。

\begin{principlebox}
先定义要比较的量，再据此规定任务、资源与分析单位。
三个百分点是历史设计参考，后续确认性研究还需按具体问题声明目标效应、精度或功效与停止规则。
试验后的规划可以指导新研究，新设计、新封存数据与历史读数各自保留清楚的证据身份。
\end{principlebox}

\falsifier{研究问题允许不同方向的答案。有效比较若支持更新有益、无实用增量或存在损害，都应按其范围报告；原假设保留在附录中供核对，新的问题表不改变历史数据和分析目标。}

\openquestion{选择下一项研究时，先指出需要识别的作用，再安排共同起点、资源范围、独立任务与分析单位。补齐运行记录、增加对照和扩大覆盖分别提供不同的信息，读者可以据此为自己的系统设计最有价值的一步。}

\section*{习题}
\begin{enumerate}[label=16-\arabic*.,ref=16-\arabic*]
 \item\label{ex:ch16-1} 为什么表~\ref{tab:ch16-power}中的区间精度规划不是检验功效设计？若要规划检验功效，还需声明哪些参数？说明为什么题数达到历史估算仍不保证可以判定。
 \item\label{ex:ch16-2} 比较研究问题P1、P3与P5：已有观察各能回答什么，下一步分别需要哪些对照或运行，哪些不能靠重复调用替代？
 \item\label{ex:ch16-3} 为 P3 设计形式与容量的交叉对照，说明任务族、检索预算和条目来源如何固定。
 \item\label{ex:ch16-4} 为 P6 定义下降量的正负号，并区分“没有检测到退化”与“有证据反对必然退化”。
 \item\label{ex:ch16-5} E03c 跨日期补跑一臂后，哪些数据变完整，哪些关于模型版本和迁移范围的问题仍未解决？
\end{enumerate}

\section*{常见误解}
\begin{itemize}
 \item \textbf{“区间跨零就是没有效果”。} 跨零不证明等价，也不自动解释为样本不足；应检查设计、精度和效应范围。
 \item \textbf{“补齐实验臂就能判命题”。} 补齐解决缺失；若对照没有隔离原主张，仍需新的设计。
 \item \textbf{“研究问题已经写成方向性结论”。} 问题用于组织比较，答案须来自对应条件下的证据；历史假设也不因改为问题而获得支持。
\end{itemize}

\section*{进一步阅读}
\begin{itemize}
  \item \cite{rsi-survey-2607}：§2的分类、§5的反馈与四类风险、§7的条件与限制，
        本章 \S\ref{sec:ch16-bounded} 与表~\ref{tab:ch16-collapse} 的骨架都在这里。
  \item \cite{self-evolving-survey-2507}：§7.1 的 $FGT_t$ / $BWT_t$ / $CPG_t$ 定义式，
        以及 §7.3.2 那句“真正的对等比较目前仍不可行”——它是本书第 3 章协议的存在理由。
  \item \cite{anatomy-memory-2602}：表 4 与表 5 给出的换主干失效与维护成本缺口，
        提示需要报告主干与实现范围，不支持所有结论只能在单一主干成立的全称断言。
  \item 附录~\ref{app:stats}：MDE、样本量反推与区间判据的推导，
        其中区分描述、支持与反证的条件。
  \item 表~\ref{tab:ch16-verdicts}：十个研究问题及下一步设计；历史假设与历次判据见附录\ref{sec:appg-history}及随书规格目录。
\end{itemize}
